\section*{SUPPLEMENTARY MATERIALS}
\subsubsection{Interval Score} uses the prediction interval, with lower and upper endpoints represented by predictive quantiles at level $\frac{\alpha}{2}$ and $1 - \frac{\alpha}{2}$. The Interval Score is defined as :
\begin{equation}
S_\alpha^{int} (l,u;y) = (u - l) +\frac{2}{\alpha}(l-y)\mathbf{1}\{y<l\} + \frac{2}{\alpha}(y-u)\mathbf{1}\{y>u\}
\end{equation}


\subsubsection{Robot grasping experiment individual experiment }

Experiment results of individual obejcts 
\begin{table}[htbp]
	\caption{Result table from grasping experiment}
	\centering
	\begin{tabular}{|l|c|c|c|c|c|c|c|c|c|c|}
		\hline
		
		 & \multicolumn{2}{c|}{\textbf{Generalized}} & \multicolumn{2}{c|}{\textbf{Gaussian}} & \multicolumn{2}{c|}{\textbf{Laplace}} & \multicolumn{2}{c|}{\textbf{Cauchy}} & \multicolumn{2}{c|}{\textbf{Evidential}} \\ \cline{2-11} 
		
		& P1 & P2 & P1 & P2 & P1 & P2 & P1 & P2 & P1 & P2 \\ \hline
		Torch & 9 & 5 & 6 & 2 & 8 & 5 & 5 & 0 & 4 & 0 \\ \hline
		Marker & 9 & 4 & 7 & 1 & 7 & 2 & 5 & 1 & 5 & 1 \\ \hline
		Mouse & 4 & 3 & 2 & 2 & 3 & 2 & 2 & 3 & 4 & 3 \\ \hline
		Brush & 5 & 4 & 4 & 3 & 6 & 4 & 3 & 1 & 3 & 0 \\ \hline
		Stapler & 6 & 7 & 1 & 4 & 3 & 5 & 1 & 3 & 1 & 2 \\ \hline
		Paste & 8 & 8 & 4 & 4 & 8 & 7 & 5 & 5 & 4 & 3 \\ \hline
		Cube & 6 & 6 & 5 & 5 & 6 & 6 & 5 & 5 & 5 & 5 \\ \hline
		Duster & 6 & 5 & 4 & 2 & 7 & 5 & 2 & 1 & 3 & 0 \\ \hline
		Box & 9 & 4 & 7 & 3 & 9 & 3 & 7 & 2 & 7 & 1 \\ \hline
		Allen & 10 & 5 & 4 & 2 & 9 & 2 & 8 & 2 & 8 & 3 \\ \hline
		
	\end{tabular}
	
\end{table}


\subsubsection*{LIMITATIONS}

The proposed loss function is not convex in general. 
To demonstrate non-convexity, consider the following scenario:
Fixed $\alpha$, Let $\alpha = 1$.
Varying $r$, Choose two points $r1 = 1$ and $r2 = -1$.
Convexity condition: A function f is convex if for all $r1$, $r2$ and $\lambda \in [0, 1]$:
\begin{equation*}
    f(\lambda x1 + (1-\lambda) x2) \leq \lambda f(x1) + (1-\lambda)L(x2)
\end{equation*}
Evaluating the loss function:
\begin{equation*}
L(r1) = 1 - log(\beta) + log(\Gamma(1/\beta))
\end{equation*}
\begin{equation*}
L(r2) = 1 - log(\beta) + log(\Gamma(1/\beta)) \\
\end{equation*}
\begin{equation*}
L(\lambda r1 + (1-\lambda)r2) = | \lambda - (1-\lambda)|^\beta - log(\beta) + log(\Gamma(1/\beta))
\end{equation*}
Checking the convexity condition:
For $\lambda$ = 0.5:
\begin{equation*}
   L(0.5r1 + 0.5r2) = 0 - log(\beta) + log(\Gamma(1/\beta)) 
\end{equation*}
\begin{multline*}
\lambda L(r1) + (1-\lambda)L(r2) \\ 
0.5(1 - log(\beta) + log(\Gamma(1/\beta))) \\ + 0.5(1 - log(\beta) + log(\Gamma(1/\beta))) = 1 - log(\beta) + log(\Gamma(1/\beta))
\end{multline*}
The convexity condition is not satisfied for all $\lambda \in [0, 1]$. Therefore, the loss function is not convex.
While the loss function is not globally convex, it might be locally convex under certain conditions. However, the general form presented here does not guarantee convexity. 
We overcome this limitation by adding a software ode by clamping the loss to the very small number as possible in the cpu. This ensures that the residual never becomes zero and thus 
